%% Equational Heyting axioms and the intermediate-logic correspondence.

\section{Intermediate logics and Heyting varieties}

\begin{annexitem}{annex:p10.intermediate.1}{thm:p10-heyting-equational}{Equations for the residual and the Boolean endpoint}
\begin{proof}
Suppose first that residuation holds.
The inequality \(a\wedge1\leq a\) gives \(a\Rightarrow a=1\).
Residuation gives \(a\wedge(a\Rightarrow b)\leq b\), hence this
meet lies below \(a\wedge b\).
Conversely, \(a\wedge b\leq a\Rightarrow b\), since its meet with
\(a\) lies below \(b\); it also lies below \(a\).
This proves the second identity.
The inequality \(a\wedge b\leq b\) similarly gives
\(b\leq a\Rightarrow b\), proving the third.
Finally, for every \(d\),
\[
\begin{aligned}
 d\leq a\Rightarrow(b\wedge c)
 &\Longleftrightarrow d\wedge a\leq b\text{ and }d\wedge a\leq c\\
 &\Longleftrightarrow d\leq(a\Rightarrow b)\wedge(a\Rightarrow c).
\end{aligned}
\]
Equality of these lower-bound conditions gives the fourth identity.

Conversely, assume the four equations.
The fourth makes \(a\Rightarrow-\) monotone: if \(b\leq c\), then
\(b=b\wedge c\), so
\(a\Rightarrow b=(a\Rightarrow b)\wedge(a\Rightarrow c)\).
If \(d\leq a\Rightarrow b\), the second equation gives
\(d\wedge a\leq a\wedge(a\Rightarrow b)=a\wedge b\leq b\).
For the reverse implication, the third equation gives
\(d\leq a\Rightarrow d\), and the first and fourth give
\[
 a\Rightarrow d
  =(a\Rightarrow a)\wedge(a\Rightarrow d)
  =a\Rightarrow(a\wedge d).
\]
Thus \(a\wedge d\leq b\) implies
\(d\leq a\Rightarrow(a\wedge d)\leq a\Rightarrow b\).
This is residuation.
Associativity, commutativity, idempotence, absorption, bounds, and
distributivity form the finite lattice-equation package; adding the
four displayed equations therefore defines \(\mathsf{HA}\) equationally.

In a Heyting algebra, \(a\wedge(a\Rightarrow0)=0\).
If also \(a\vee(a\Rightarrow0)=1\), every element has a complement,
so the bounded distributive lattice is Boolean.
Conversely, in a Boolean algebra
\[
 d\wedge a\leq b\quad\Longleftrightarrow\quad d\leq\neg a\vee b.
\]
For the forward direction decompose
\(d=(d\wedge a)\vee(d\wedge\neg a)\); the reverse follows by
meeting with \(a\) and distributing.
Uniqueness of the residual gives \(a\Rightarrow b=\neg a\vee b\).
The additional identity follows.  These calculations include the
one-element algebra.
\end{proof}
\end{annexitem}

\begin{annexitem}{annex:p10.intermediate.2}{thm:p10-intermediate-algebra-completeness}{The quotient for an arbitrary intermediate logic}
\begin{proof}
\paragraph{Congruence and order.}
The implication rules of IPC prove reflexivity and the transitivity
schema
\[
 (\alpha\to\beta)\to((\beta\to\gamma)\to(\alpha\to\gamma)).
\]
IPC also proves conjunction introduction, both projections, and the
schemas expressing monotonicity of \(\wedge,\vee\) in their arguments
and antitonicity/monotonicity of implication.  The latter is
\[
 ((\alpha'\to\alpha)\wedge(\beta\to\beta'))
       \to((\alpha\to\beta)\to(\alpha'\to\beta')).
\]
For its derivation, assume the conjunction and \(\alpha\to\beta\),
then assume \(\alpha'\); implication elimination yields successively
\(\alpha,\beta,\beta'\).  Discharge the assumptions by implication
introduction.  The conjunction monotonicity schema follows by taking
the two projections, applying the two assumed implications, and
reintroducing conjunction.  For disjunction, apply the appropriate
implication in each disjunctive case and introduce the target disjunction.
These are instances of the propositional LJ rules used in
Annex~\ref{annex:p10.heyting.2}.

All these schemas belong to \(L\).  Applying them by modus ponens
shows that theorem implication is a preorder, that its symmetric part
is \(\equiv_L\), and that equivalent arguments give equivalent results
for each connective.  Thus \(\equiv_L\) is a congruence and the
displayed quotient order is well-defined and antisymmetric.
IPC's bottom and top laws, conjunction projections/introduction, and
disjunction introductions/elimination give the bounds, meets, and
joins in this order.  Its distributive formulas give distributivity.

IPC proves the currying equivalence
\[
 ((\chi\wedge\alpha)\to\beta)
       \leftrightarrow(\chi\to(\alpha\to\beta)).
\]
In each direction, assume the left or right implication and then its
needed arguments; conjunction and implication elimination give
\(\beta\), and implication introduction discharges them.
Since both directions are in \(L\), modus ponens gives
\[
 [\chi]_L\wedge[\alpha]_L\leq[\beta]_L
       \quad\Longleftrightarrow\quad
 [\chi]_L\leq[\alpha\to\beta]_L.
\]
Every quotient element has a representative, so this proves the
residual law at every triple.  Hence \(H_L\) is a Heyting algebra.
No substitution on a premise-dependent line has been used.

\paragraph{Theoremhood and arbitrary valuations.}
The IPC theorems \(\lambda\to\top\) and
\(\lambda\to(\top\to\lambda)\), together with conjunction
introduction as an IPC theorem schema, give
\[
 [\lambda]_L=1\quad\Longleftrightarrow\quad\lambda\in L.
\]
Indeed, a theorem \(\lambda\) gives both implications between
\(\lambda\) and \(\top\); conversely their equivalence and the
theorem \(\top\) give \(\lambda\) by modus ponens.
The canonical valuation \(\nu(p)=[p]_L\) satisfies
\(\llbracket\varphi\rrbracket_\nu=[\varphi]_L\), by formula induction.
Because \(\bot\notin\mathsf{CPL}\) and \(L\subseteq\mathsf{CPL}\),
the quotient is nontrivial.

Fix \(\lambda\in L\) and any valuation \(v:X\to H_L\).
Choose representative formulas \(\rho_p\) for the values \(v(p)\)
only for the finitely many variables in \(\lambda\).
Let \(\sigma(p)=\rho_p\) there and \(\sigma(p)=p\) elsewhere.
Formula induction gives
\[
 \llbracket\lambda\rrbracket_v=[\sigma(\lambda)]_L.
\]
Uniform substitution is a theorem closure rule for \(L\), so
\(\sigma(\lambda)\in L\), and the displayed value is \(1\).
Thus \(H_L\in\mathsf V(L)\).  This argument requires only finite
choices of representatives, not a simultaneous choice for the quotient.
If \(\varphi\notin L\), its canonical value is not \(1\), so
\(H_L\) refutes it.  Conversely every member of \(L\) is valid in
\(\mathsf V(L)\) by definition.  This proves all claimed equivalences.
\end{proof}
\end{annexitem}

\begin{annexitem}{annex:p10.intermediate.3}{thm:p10-intermediate-variety-correspondence}{The inverse assignments and their lattice laws}
\begin{proof}
\paragraph{Equations and formulas.}
For a \(\tau_H\)-term \(s\), let \(s^\flat\) be the formula obtained
by replacing \(0,1,\Rightarrow\) with \(\bot,\top,\to\), leaving
variables, meets, and joins unchanged.
Formula induction identifies its value with the term value.
Residuation gives
\(a\Rightarrow b=1\) exactly when \(a\leq b\).
Consequently, for any Heyting algebra and valuation,
\[
 s=t\quad\Longleftrightarrow\quad
 (s\Rightarrow t)\wedge(t\Rightarrow s)=1.
\]
Thus the identity \(s\approx t\) is equivalent to validity of
\(s^\flat\leftrightarrow t^\flat\), and validity of a formula
\(\varphi\) is the identity \(\varphi^{\rm alg}\approx1\), where
\(\varphi^{\rm alg}\) replaces its connectives by the algebra operations.
The equivalences quantify over all valuations, not one fixed assignment.
They include ground equations because \(\top\) represents the constant \(1\).

It follows that \(\mathsf V(L)\) is a subvariety of the equational
class \(\mathsf{HA}\), axiomatized there by
\(\lambda^{\rm alg}\approx1\) for \(\lambda\in L\).

\paragraph{The endpoints.}
Theorem~\ref{thm:p10-heyting-completeness} gives both
\(\mathsf V(\mathsf{IPC})=\mathsf{HA}\) and
\(\operatorname{Log}(\mathsf{HA})=\mathsf{IPC}\).
In a Boolean Heyting algebra the implication is \(\neg a\vee b\),
by the preceding equational theorem.
Every classical tautology is valid in every such algebra.
For an explicit verification, if its variables have values
\(a_1,\ldots,a_m\), form the Boolean cells
\[
 c_\epsilon=\bigwedge_{i=1}^m a_i^{\epsilon_i},
 \qquad a_i^1=a_i,\quad a_i^0=\neg a_i.
\]
Their join is \(1\).  Induction on a formula shows that its value,
met with \(c_\epsilon\), is \(c_\epsilon\) or \(0\) according to
its classical value under \(\epsilon\); implication uses
\(\neg a\vee b\).  A tautology therefore has value \(1\).
This includes the zero-variable case, with the sole cell \(1\).
Conversely the two-element Boolean algebra refutes every formula
outside CPL, by classical propositional completeness from Part~II.
Hence \(\operatorname{Log}(\mathsf{BA})=\mathsf{CPL}\).

A Heyting algebra validating CPL validates \(p\vee\neg p\), so
\(a\vee(a\Rightarrow0)=1\) for every element \(a\); the equational
theorem makes it Boolean.  The reverse inclusion follows from the
tautology calculation.  Thus \(\mathsf V(\mathsf{CPL})=\mathsf{BA}\).
In particular, \(\mathsf{BA}\subseteq\mathsf V(L)\subseteq\mathsf{HA}\).

If \(\mathsf{BA}\subseteq\mathcal V\subseteq\mathsf{HA}\), these
endpoint identities imply
\(\mathsf{IPC}\subseteq\operatorname{Log}(\mathcal V)\subseteq\mathsf{CPL}\).
Modus ponens preserves value \(1\), since \(a=1\) and
\(a\Rightarrow b=1\) force \(b=1\).
Uniform substitution preserves validity by composing a valuation
with formula evaluation.  Thus \(\operatorname{Log}(\mathcal V)\)
is an intermediate logic.

\paragraph{Both inverse equalities.}
The quotient completeness theorem gives
\(\operatorname{Log}(\mathsf V(L))=L\).
For the other equality, let \(\mathcal V\) be axiomatized within
\(\mathsf{HA}\) by identities \(E\).
Every member of \(\mathcal V\) validates its logic, giving
\(\mathcal V\subseteq\mathsf V(\operatorname{Log}(\mathcal V))\).
For \(s\approx t\in E\), the equation--formula translation puts
\(s^\flat\leftrightarrow t^\flat\) in \(\operatorname{Log}(\mathcal V)\).
Any member of \(\mathsf V(\operatorname{Log}(\mathcal V))\) therefore
satisfies every identity in \(E\), and belongs to \(\mathcal V\).
This proves the reverse inclusion.

\paragraph{Order and complete lattice operations.}
Adding theorems can only remove validating algebras, and adding
algebras can only remove universally valid formulas.
Thus both assignments reverse inclusion.  The inverse equalities
make them order-reversing bijections.

Intersections of nonempty families of intermediate logics preserve
the endpoint inclusions and both closure rules.  The least closed
extension of a union exists by intersecting its intermediate upper
bounds; CPL is one such upper bound.  This gives all logic meets
and joins, with the stated empty-family endpoints.
Intersections of subvarieties are equationally definable by unions
of their canonical identity theories: for each subvariety take all
identities valid in it.  This avoids selecting one axiomatization
for each member of an infinite family.  A least subvariety containing
\(\mathsf{BA}\) and a given union exists by taking all identities
valid on that class and then all their Heyting models.
Indeed, each upper-bound subvariety has a defining identity set
among these identities, so the resulting class is below every upper bound.
This gives the variety lattice and its stated empty-family endpoints.
An order-reversing bijection exchanges greatest lower bounds with
least upper bounds: translate the lower-bound inequalities by either
assignment and use their inverse equalities.  Hence it exchanges
all meets and joins, including the empty ones.
All arguments here use ZF and finite representative choices;
no prime-filter or Kripke representation is required.
\end{proof}
\end{annexitem}
